% ==============================================================================
% MÓDULO CAPÍTULO II: CARACTERÍSTICAS DEMOGRÁFICAS Y ESTRUCTURA POBLACIONAL
% Archivo: cap_2/2.2_demografia_poblacion.tex
% Enfoque: 100% Calidad de Atención (Donabedian y Watson)
% ==============================================================================
\subsection{Características demográficas y estructura poblacional}
\label{sec:2.2_demografia}

La jurisdicción sanitaria asignada al Centro de Salud Belén abarca una población diana proyectada de 18,450 habitantes según el Padrón Poblacional Oficial de la Red de Salud Huamanga (2025--2026). La pirámide poblacional se caracteriza por una base ensanchada en los primeros decenios de vida y una proporción creciente de personas adultas mayores, configurando un doble desafío asistencial de perfil materno-infantil y crónico-degenerativo. La distribución por etapas de vida es la siguiente:
\begin{itemize}
    \item \textbf{Población Infantil (menores de 5 años):} Representa el 12.4\% (2,288 niños), constituyendo el segmento con mayor intensidad de demanda en los consultorios preventivo-promocionales de Crecimiento y Desarrollo (CRED), tamizaje y tratamiento de anemia ferropénica, suplementación nutricional e inmunizaciones del esquema nacional regular.
    \item \textbf{Población Escolar y Adolescente (5 a 17 años):} Abarca el 19.6\% (3,616 personas), demandantes de intervenciones de salud bucal, tamizaje visual, salud del escolar y orientación en salud sexual y reproductiva.
    \item \textbf{Población Joven y Adulta (18 a 59 años):} Constituye el 52.8\% (9,742 personas), con un marcado predominio de mujeres en edad fértil (MEF) que acuden para control prenatal, despistaje de cáncer de cuello uterino, planificación familiar y atenciones médicas por patologías agudas.
    \item \textbf{Población Adulta Mayor (60 años a más):} Concentra el 15.2\% (2,804 personas), un grupo en acelerado envejecimiento que presenta alta prevalencia de hipertensión arterial, diabetes mellitus tipo 2, artrosis y secuelas osteomusculares, requiriendo un acompañamiento asistencial continuo, trato digno y calidez humana preferencial.
\end{itemize}

Desde la mirada fenomenológica, resalta el fenómeno de la \textbf{feminización de la demanda matutina}: más del 70\% de los usuarios que concurren a las salas de espera son mujeres. Se trata de madres de familia que acuden al cuidado y vigilancia de la salud de sus hijos menores, mujeres gestantes en controles prenatales y mujeres cuidadoras que asisten a adultos mayores dependientes. Esta dinámica impone sobre la mujer andina una sobrecarga cotidiana de roles, donde la calidad del cuidado recibida en enfermería y medicina —medida en paciencia explicativa, respeto a su dignidad y calidez en el trato— resulta decisiva para afianzar su confianza en los servicios del Estado.
